\SourceLeaf{65}
\begin{center}
\textit{SECHSTER BEWEIS DES FUNDAMENTALTHEOREMS}\par
\textit{IN DER THEORIE DER QUADRATISCHEN RESTE.}
\end{center}

\Article{1}

\textsc{Satz.}\quad \textit{Es bezeichne $p$ eine Primzahl (positiv und ungerade), $n$ eine positive ganze, durch $p$ nicht teilbare Zahl und $x$ eine unbestimmte Größe. Dann ist die Funktion}

\[
\begin{aligned}
1+x^n+x^{2n}+x^{3n}+\text{usw.}+x^{np-n}
\end{aligned}
\]

\textit{teilbar durch}

\[
\begin{aligned}
1+x+xx+x^3+\text{usw.}+x^{p-1}
\end{aligned}
\]

{\rightskip=0pt plus 1fil \textsc{Beweis.}\quad Man nehme eine positive ganze Zahl $g$ so, dass $gn\equiv1\quad(\text{mod. }p)$ gilt, und setze $gn=1+hp$.\quad Dann gilt\par}

{\small
\[
\begin{aligned}
\frac{1+x^n+x^{2n}+x^{3n}+\text{usw.}\,x^{np-n}}{1+x+xx+x^3+\text{usw.}\,x^{p-1}} &=\frac{(1-x^{np})(1-x)}{(1-x^n)(1-x^p)}\\&=\frac{(1-x^{np})(1-x^{gn}-x+x^{hp+1})}{(1-x^n)(1-x^p)}\\
&=\frac{1-x^{np}}{1-x^p}\cdot\frac{1-x^{gn}}{1-x^n}-\frac{x(1-x^{np})}{1-x^n}\cdot\frac{1-x^{hp}}{1-x^p}
\end{aligned}
\]
}

und dieser Ausdruck ist offenbar ein Polynom.\quad Q. E. D.

Jedes Polynom in $x$, das durch $\frac{1-x^{np}}{1-x^n}$ teilbar ist, ist daher auch durch $\frac{1-x^p}{1-x}$ teilbar.

\Article{2}

Es bezeichne $\alpha$ eine positive primitive Wurzel modulo $p$; das heißt, $\alpha$ sei eine positive ganze Zahl, deren Potenzen $1,\alpha,\alpha\alpha,\alpha^3\dotsfive\alpha^{p-2}$ modulo $p$ als kleinste positive Reste ohne Rücksicht auf die Reihenfolge genau die Zahlen $1,2,3,4\dotsfive p-1$ haben.\quad Bezeichnen wir ferner mit $fx$ die Funktion

\[
\begin{aligned}
x+x^\alpha+x^{\alpha\alpha}+x^{\alpha^3}+\text{usw.}+x^{\alpha^{p-2}}+1
\end{aligned}
\]

so ist offenbar $fx-1- x-xx-x^3-\text{usw.}-x^{p-1}$ durch $1-x^p$ teilbar und damit erst recht durch $\frac{1-x^p}{1-x}=1+x+xx+x^3+\text{usw.}+x^{p-1}$. Durch diese letzte Funktion ist daher auch $fx$ selbst teilbar.\quad Da $x$ eine unbestimmte Größe darstellt, folgt daraus, dass auch $f(x^n)$ durch $\frac{1-x^{np}}{1-x^n}$ und somit nach dem vorhergehenden Artikel auch durch $\frac{1-x^p}{1-x}$ teilbar ist, sofern $n$ eine ganze, durch $p$ nicht teilbare Zahl ist. Ist dagegen $n$ eine ganze, durch $p$ teilbare Zahl, so sind die einzelnen Glieder der Funktion $f(x^n)$, jeweils um\SourceLeaf{66} eins vermindert, durch $1-x^p$ teilbar. Daher ist in diesem Fall auch $f(x^n)-p$ durch $1-x^p$ und somit auch durch $\frac{1-x^p}{1-x}$ teilbar.

\Article{3}

\textsc{Satz.}\quad \textit{Setzt man}

\[
\begin{aligned}
x-x^\alpha+x^{\alpha\alpha}-x^{\alpha^3}+x^{\alpha^4}-\text{usw.}-x^{\alpha^{p-2}}=\xi
\end{aligned}
\]

\textit{so ist $\xi\xi\mp p$ durch $\frac{1-x^p}{1-x}$ teilbar. Dabei ist das obere Vorzeichen zu nehmen, wenn $p$ die Form $4k+1$ hat, das untere, wenn $p$ die Form $4k+3$ hat.}

\textsc{Beweis.}\quad Man sieht leicht, dass von den folgenden $p-1$ Funktionen

\[
\begin{aligned}
&+x\xi-xx+x^{\alpha+1}-x^{\alpha\alpha+1}\\&\qquad+\text{usw.}+x^{\alpha^{p-2}+1}\\
&-x^\alpha\xi-x^{2\alpha}+x^{\alpha\alpha+\alpha}-x^{\alpha^3+\alpha}\\&\qquad+\text{usw.}+x^{\alpha^{p-1}+\alpha}\\
&+x^{\alpha\alpha}\xi-x^{2\alpha\alpha}+x^{\alpha^3+\alpha\alpha}-x^{\alpha^4+\alpha\alpha}\\&\qquad+\text{usw.}+x^{\alpha^p+\alpha\alpha}\\
&-x^{\alpha^3}\xi-x^{2\alpha^3}+x^{\alpha^4+\alpha^3}-x^{\alpha^5+\alpha^3}\\&\qquad+\text{usw.}+x^{\alpha^{p+1}+\alpha^3}
\end{aligned}
\]

usw. bis zu

\[
\begin{aligned}
-x^{\alpha^{p-2}}\xi-x^{2\alpha^{p-2}}+x^{\alpha^{p-1}+\alpha^{p-2}}-x^{\alpha^p+\alpha^{p-2}}\\&\qquad+\text{usw.}+x^{\alpha^{2p-4}+\alpha^{p-2}}
\end{aligned}
\]

die erste $=0$ ist, jede der übrigen dagegen durch $1-x^p$ teilbar. Daher ist auch ihre Summe durch $1-x^p$ teilbar; sie ergibt sich zu

\[
\begin{aligned}
&=\xi\xi-(f(xx)-1)+(f(x^{\alpha+1})-1)-(f(x^{\alpha\alpha+1})-1)\\&\qquad+(f(x^{\alpha^3+1})-1)-\text{usw.}\\
&\qquad+(f(x^{\alpha^{p-2}+1})-1)\\
&=\xi\xi-f(xx)+f(x^{\alpha+1})-f(x^{\alpha\alpha+1})\\&\qquad+f(x^{\alpha^3+1})-\text{usw.}+f(x^{\alpha^{p-2}+1})=\Omega
\end{aligned}
\]

\Needspace{13\baselineskip}
Damit ist dieser Ausdruck $\Omega$ auch durch $\frac{1-x^p}{1-x}$ teilbar. Unter den Exponenten $2,\ \alpha+1,\ \alpha\alpha+1,\ \alpha^3+1\dotsfive\alpha^{p-2}+1$ ist nun nur einer durch $p$ teilbar, nämlich $\alpha^{\frac12(p-1)}+1$. Nach dem vorhergehenden Artikel sind daher die folgenden einzelnen Glieder des Ausdrucks $\Omega$

\[
\begin{aligned}
f(xx),\quad f(x^{\alpha+1}),\quad f(x^{\alpha\alpha+1}),\quad (fx^{\alpha^3+1})\quad\text{usw.}
\end{aligned}
\]

mit alleiniger Ausnahme des Gliedes $f(x^{\alpha^{\frac12(p-1)}+1})$ durch $\frac{1-x^p}{1-x}$ teilbar. Diese Glieder dürfen daher weggelassen werden, so dass die folgende Funktion weiterhin durch $\frac{1-x^p}{1-x}$ teilbar bleibt:

\[
\begin{aligned}
\xi\xi\mp f(x^{\alpha^{\frac12(p-1)}+1})
\end{aligned}
\]

\SourceLeaf{67} Dabei gilt das obere oder untere Vorzeichen, je nachdem, ob $p$ die Form $4k+1$ oder die Form $4k+3$ hat. Da außerdem $f(x^{\alpha^{\frac12(p-1)}+1})-p$ durch $\frac{1-x^p}{1-x}$ teilbar ist, ist auch $\xi\xi\mp p$ durch $\frac{1-x^p}{1-x}$ teilbar. Q. E. D.

Damit das doppelte Vorzeichen keine Zweideutigkeit hervorruft, bezeichnen wir mit $\varepsilon$ die Zahl $+1$ oder $-1$, je nachdem, ob $p$ die Form $4k+1$ oder $4k+3$ hat. Damit ist $\frac{(1-x)(\xi\xi-\varepsilon p)}{1-x^p}$ ein Polynom in $x$, das wir mit $Z$ bezeichnen.

\Article{4}

Es sei $q$ eine positive ungerade Zahl, also $\frac12(q-1)$ eine ganze Zahl. Dann ist $(\xi\xi)^{\frac12(q-1)}-(\varepsilon p)^{\frac12(q-1)}$ durch $\xi\xi-\varepsilon p$ und somit auch durch $\frac{1-x^p}{1-x}$ teilbar. Wir setzen $\varepsilon^{\frac12(q-1)}=\delta$ und

\[
\begin{aligned}
\xi^{q-1}-\delta p^{\frac12(q-1)}=\frac{1-x^p}{1-x}\cdot Y
\end{aligned}
\]

Dann ist $Y$ ein Polynom in $x$, und es gilt $\delta=\mathord{.}+1$, wenn für eine der Zahlen $p,q$ oder auch für beide die Form $4k+1$ vorliegt; dagegen gilt $\delta=-1$, wenn beide $p,q$ die Form $4k+3$ haben.

\newpage
\Article{5}

Nehmen wir nun an, auch $q$ sei eine Primzahl, die von $p$ verschieden ist. Nach dem in den \textit{Disquisitiones Arithmeticae}, Art. 51, bewiesenen Satz ist offenbar

\[
\begin{aligned}
\mathord{.}\quad\xi^q-(x^q-x^{q\alpha}+x^{q\alpha\alpha}-x^{q\alpha^3}+\text{usw.}-x^{q\alpha^{p-2}})
\end{aligned}
\]

durch $q$ teilbar oder von der Form $qX$, wobei $X$ ein Polynom in $x$ ist, dessen numerische Koeffizienten ebenfalls ganze Zahlen sind (was auch für die übrigen hier vorkommenden Polynome $Z,Y,W$ vorauszusetzen ist). Für den Modul $p$ und die primitive Wurzel $\alpha$ bezeichnen wir den Index der Zahl $q$ mit $\mu$; das heißt, es gelte $q\equiv\alpha^\mu\ (\mathrm{mod.}\ p)$. Die Zahlen $q,\ q\alpha,\ q\alpha\alpha,\ q\alpha^3\dotsfive q\alpha^{p-2}$ sind daher modulo $p$ jeweils den Zahlen $\alpha^\mu,\allowbreak\ \alpha^{\mu+1},\allowbreak\ \alpha^{\mu+2}\dotsfour\alpha^{p-2},\allowbreak\ 1,\allowbreak\ \alpha,\allowbreak\ \alpha\alpha\dotsfour\alpha^{\mu-1}$ kongruent, also sind

\[
\begin{aligned}
&x^q-x^{\alpha^\mu}\\
&x^{q\alpha}-x^{\alpha^{1+1}}\\
&x^{q\alpha\alpha}-x^{\alpha^{\mu+2}}\\
&x^{q\alpha^3}-x^{\alpha^{\mu+3}}\\
&\hspace{9mm}\begin{gathered}.\\.\\.\\.\end{gathered}
\end{aligned}
\]

\SourceLeaf{68}
\[
\begin{aligned}
&x^{q\alpha^{p-\mu-2}}-x^{\alpha^{p-2}}\\
&x^{q\alpha^{p-\mu-1}}-x\\
&x^{q\alpha^{p-\mu}}-x^\alpha\\
&x^{q\alpha^{p-\mu+1}}-x^{\alpha\alpha}\\
&\hspace{9mm}\begin{gathered}.\\.\\.\\.\end{gathered}\\
&x^{q\alpha^{p-2}}-x^{\alpha^{\mu-1}}
\end{aligned}
\]

\Needspace{14\baselineskip}
durch $1-x^p$ teilbar. Nehmen wir diese Größen abwechselnd mit positivem und negativem Vorzeichen und addieren sie, so sieht man, dass $1-x^p$ die Funktion

\[
\begin{aligned}
x^q-x^{q\alpha}+x^{q\alpha\alpha}-x^{q\alpha^3}+\text{usw.}-x^{q\alpha^{p-2}}\mp\xi
\end{aligned}
\]

teilt. Dabei gilt das obere oder untere Vorzeichen, je nachdem, ob $\mu$ gerade oder ungerade ist, das heißt, ob $q$ ein quadratischer Rest oder Nichtrest von $p$ ist. Wir setzen daher

\[
\begin{aligned}
x^q-x^{q\alpha}+x^{q\alpha\alpha}- x^{q\alpha^3}+\text{usw.}-x^{q\alpha^{p-2}}-\gamma\xi\\&=(1-x^p)W
\end{aligned}
\]

mit $\gamma=+1$ oder $\gamma=-1$, je nachdem, ob $q$ ein quadratischer Rest oder Nichtrest von $p$ ist; dabei ist $W$ offenbar ein Polynom.

\Needspace{13\baselineskip}
\Article{6}

Nach diesen Vorbereitungen leiten wir aus der Verbindung der vorhergehenden Gleichungen die folgende her:

\[
\begin{aligned}
q\xi X=\varepsilon p(\delta p^{\frac12(q-1)}-\gamma)\\&\qquad+\frac{1-x^p}{1-x}\cdot(Z(\delta p^{\frac12(q-1)}-\gamma)\\&\qquad+Y\xi\xi-W\xi(1-x))
\end{aligned}
\]

Nehmen wir an, dass die Division der Funktion $\xi X$ durch

\[
\begin{aligned}
x^{p-1}+x^{p-2}+x^{p-3}+\text{usw.}+x+1
\end{aligned}
\]

den Quotienten $U$ mit dem Rest $T$ ergibt, oder dass gilt:

\[
\begin{aligned}
\xi X=\frac{1-x^p}{1-x}\cdot U+T
\end{aligned}
\]

Dabei sind $U,T$ Polynome, deren numerische Koeffizienten ebenfalls ganze Zahlen sind, und $T$ hat einen sicher niedrigeren Grad als der Divisor. Dann gilt

\[
\begin{aligned}
qT-\varepsilon p(\delta p^{\frac12(q-1)}-\gamma)\\&=\frac{1-x^p}{1-x}\cdot(Z(\delta p^{\frac12(q-1)}-\gamma)\\&\qquad+Y\xi\xi-W\xi(1-x)-qU)
\end{aligned}
\]

\Needspace{8\baselineskip}
Diese Gleichung kann offenbar nur bestehen, wenn sowohl die linke als auch die rechte Seite für sich verschwindet. Daher ist $\varepsilon p(\delta p^{\frac12(q-1)}-\gamma)$ durch $q$ teilbar\SourceLeaf{69}, ebenso auch $\delta p^{\frac12(q-1)}-\gamma$; somit ist wegen $\delta\delta=1$ auch die Zahl $p^{\frac12(q-1)}-\gamma\delta$ durch $q$ teilbar.

Bezeichnet nun $\beta$ die positiv oder negativ genommene Einheit, je nachdem, ob $p$ ein quadratischer Rest oder Nichtrest der Zahl $q$ ist, so ist $p^{\frac12(q-1)}-\beta$ durch $q$ teilbar und damit auch $\beta-\gamma\delta$. Dies ist nur möglich, wenn $\beta=\gamma\delta$ gilt. Daraus folgt unmittelbar das Fundamentaltheorem. Nämlich:

I. Wenn die Zahlen $p,q$ beide die Form $4k+1$ haben oder nur eine von ihnen, so dass $\delta=+1$ gilt, ist $\beta=\gamma$. Daher ist entweder zugleich $q$ ein quadratischer Rest von $p$ und $p$ ein quadratischer Rest von $q$, oder zugleich $q$ ein Nichtrest von $p$ und $p$ ein Nichtrest von $q$.

II. Wenn beide $p,q$ die Form $4k+3$ haben, so dass $\delta=-1$ gilt, ist $\beta=-\gamma$. Daher ist entweder zugleich $q$ ein quadratischer Rest von $p$ und $p$ ein Nichtrest von $q$, oder zugleich $q$ ein Nichtrest von $p$ und $p$ ein Rest von $q$.

Q. E. D.
